\begin{enumerate}
  \item Squad Viva GeoAnalytic a.k.a VGA---Rico, Agus, Zaky, Louis, Ilyas, Dhika, Dhanar, dan Tito---atas segala gendongannya agar penulis dapat mencapai rank Mythical Immortal sebagai dukungan moral selama pengerjaan Tugas Akhir ini.
  \item Teman-teman Rambaklicious---Adinda, Elvina, Hafidz, Noni, Khamim, dan Zalfa---atas segala dukungan dan kebersamaannya selama Tugas Akhir ini.
  \item Teman-teman coc man---Alfian, Chaca, Ananda, Rizka, Fira, Melinda, dan Jannah---selaku teman seperjuangan saat kelas aksel di MAN Surabaya yang telah memberikan dukungan dan motivasi selama perkuliahan.
  \item Teman-teman Asisten Laboratorium Aljabar dan Analisis---Haidar, Fajar, Vinsen, dkk---yang senantiasa menemani penulis saat menginap di laboratorium untuk menyelesaikan Tugas Akhir ini.
  \item Mas Vian, Mas Ilham, dan Mas Hisbu atas \textit{sharing} ilmu dan pengalaman terkait perkuliahan dan penelitian. Pastinya juga terima kasih atas segala motivasinya terlebih untuk mendorong penulis mengikuti lomba-lomba Matematika.
  \item Junika dan Huda, sebagai partner lomba dari FMC UM hingga ONMIPA Bidang Matematika. Semoga kalian berdua selalu sukses dan bahagia dalam setiap langkahnya.
  \item Teman-teman Balerina Lontar---Ella, Sultan, Reyhan, Juan,dan Reza---yang telah saling memberikan \textit{vibes} positif disaat masa-masa sulit dalam pelaksanaan kegiatan HIMATIKA.
  \item Delen, Isa, dan Narendra selaku teman sekelas di SMPN 35 Surabaya yang telah senantiasa mengajak penulis untuk nongkrong dan berbagi cerita di masa-masa sulit dalam perkuliahan.
  \item Question Maker OMITS 16th dan Question Maker 17th, yang telah memberikan pengalaman berharga dalam menyusun soal-soal OMITS.
  \item Bang Farrel (Mao), yang selalu memberikan rekomendasi terkait barang elektronik apapun hingga ikut menemani penulis saat membeli barang elektronik secara COD.
  \item Bu Alvida Mustika Rukmi, S.Si, M.Si, selaku dosen yang banyak memberikan amanah kepada penulis dalam berbagai kegiatan akademik maupun non-akademik, sehingga penulis dapat belajar banyak hal dari beliau.
  \item Mark Zuckerberg, selaku pendiri Facebook, yang telah menciptakan platform media sosial untuk menjadi sarana komunikasi dan hiburan bagi penulis selama pengerjaan Tugas Akhir ini.
  \item Fajrul Fx, Wielino, DIMSK, dan Bangskot selaku kreator konten di YouTube yang telah memberikan ilmu dan hiburan bagi penulis selama pengerjaan Tugas Akhir ini.
  \item AliA, Yorushika, Rokudenashi, Mrs. Green Apple, Chico with Honeyworks, tuki., Ryokuoushoku Shakai, SPYAIR, Yuuri, Yuika, Ado, dan berbagai musisi lainnya yang telah menciptakan lagu-lagu untuk menemani penulis saat mengerjakan Tugas Akhir ini.
  \item Hoyoverse, Moonton, Nexon dan Supercell yang telah menciptakan game-game seru untuk menemani penulis saat mengerjakan Tugas Akhir ini.
\end{enumerate}